% TOP_PROBLEM: {"id": 20000889, "title": "A holonomy reduction to equiangular lines", "category_id": 2, "category": {"id": 2, "name": "combinatorics", "display_name": "Combinatorics", "description": "Counting problems, graph theory, discrete structures.", "slug": "combinatorics", "order_index": 2, "created_at": "2026-07-31T15:26:25.671Z"}, "status": "partially_solved", "problem_number": "AIM-COMBINATORICS-0014", "set_id": 14}
% FIRST_POSTED: 2026-10-01
% ENTRY_KIND: statement_only
% UnsolvedMath ID 20000889; source catalogue code AIM-COMBINATORICS-0014
% !TeX program = lualatex
% Definition and sources only; this file is not an LLM solution attempt.
\documentclass[11pt,a4paper]{article}
\usepackage[margin=25mm]{geometry}
\usepackage{fontspec}
\setmainfont{lmroman10-regular.otf}[BoldFont=lmroman10-bold.otf,
 ItalicFont=lmroman10-italic.otf,BoldItalicFont=lmroman10-bolditalic.otf]
\usepackage{amsmath,amssymb,mathtools}
\usepackage{unicode-math}
\setmathfont{latinmodern-math.otf}
\AtBeginDocument{\let\setminus\smallsetminus}
\usepackage{xurl,hyperref}
\hypersetup{colorlinks=true,urlcolor=blue}
\setcounter{secnumdepth}{0}
\setlength{\parindent}{0pt}
\setlength{\parskip}{5pt}
\setlength{\emergencystretch}{3em}
\begin{document}
\section*{A holonomy reduction to equiangular lines}\label{20000889}
{\small UnsolvedMath ID 20000889 · AIM-COMBINATORICS-0014 · Combinatorics · AIM Workshop Problem Lists}
\subsection{Definitions and mathematical statement}
Let $1\le k\le d$ and $0\le\alpha\le1$. For a linear subspace
$V\subseteq\mathbb R^d$, let $P_V$ be orthogonal projection onto $V$.
Two $k$-dimensional subspaces $U,V$ are isoclinic with parameter $\alpha$
if $\|P_Vu\|=\alpha$ for every unit vector $u\in U$. Equivalently,
\[
                P_UP_VP_U=\alpha^2P_U.
\]
For equal-dimensional subspaces this condition is symmetric. The
parameter used in the workshop is a projection length; the corresponding
principal angle is $\arccos\alpha$.
Let $N_\alpha^k(d)$ be the largest cardinality of a family of distinct
$k$-dimensional subspaces of $\mathbb R^d$ in which every pair satisfies
this condition with the same $\alpha$.

The two proposed bounds are:
\begin{enumerate}
\item For every fixed $k$ and $\alpha$, there is a finite constant
$C_{k,\alpha}$ such that $N_\alpha^k(d)\le C_{k,\alpha}d$ for all $d\ge k$.
\item For every fixed $k$, there is a constant $C_k$ such that for each
fixed $\alpha$,
$\limsup_{d\to\infty}N_\alpha^k(d)/d\le C_k$.
\end{enumerate}
The second assertion expresses $O_k(d)+o_{\alpha,k}(d)$: its leading
coefficient is independent of $\alpha$, although the dimension from which
an approximate bound holds may depend on $\alpha$. The substantive range
is $0<\alpha<1$; at $\alpha=0$ the subspaces are mutually orthogonal,
and at $\alpha=1$ two distinct subspaces cannot satisfy the condition.
\subsection{Short English statement}
Must an equi-isoclinic family of fixed-dimensional real subspaces grow only linearly with the ambient dimension, with a leading constant independent of the common angle?
\subsection{Sources}
\begin{itemize}
\item[S1] ULAM.AI and the UnsolvedMath Contributors, supplied problems.json, record 20000889 (AIM-COMBINATORICS-0014), AIM Workshop Problem Lists. Catalogue identity and source transcription; CC BY 4.0.
\url{https://huggingface.co/datasets/ulamai/UnsolvedMath}.
\item[S2] American Institute of Mathematics, High-dimensional phenomena in discrete analysis, item 4.1. Original workshop question underlying AIM-COMBINATORICS-0014.
\url{http://aimpl.org/highdimdiscrete/4/}.
\end{itemize}
\end{document}
